\subsection*{把未知系数藏进面积比}
目标写作 $b=xu+yv$，其中 $u,v$ 是矩阵的列。除了逐坐标解方程，还能寻找一个操作，只留下 $x$ 而消去 $y$。以 $b,v$ 为边计算有向面积：$yv$ 与 $v$ 平行，不贡献面积；剩余的 $xu$ 使原面积乘以 $x$。

用行列式语言写出刚才的几何论证：
\[\det[b\ v]=\det[xu+yv\ v]=x\det[u\ v].\]
原面积不为零时，除过去便得到 $x$；类似地，以 $u,b$ 为列得到 $y$。“换列后求行列式”其实是一种消去另一个未知数的测量。

\subsection*{计算之后仍要返回原问题}
取 $u=(2,1),v=(1,3),b=(5,5)$。原面积为 $5$，替换第一列后为 $10$，故第一系数为 $2$；替换第二列后为 $5$，另一系数为 $1$。最后代回 $2u+v=(5,5)$，检查列序与符号。

若原面积为零，除法失效反映问题已经改变：两列只张成一条线，目标可能不可达，也可能有多套表示。应回到列空间与核，而不能把零分母约掉。几何解释恰好说明了公式条件为什么必要。

这套公式适合建立唯一解与有向面积的联系，并不意味着大规模数值求解都应采用它；效率与误差稳定性需要另外研究。练习时先画列向量，再把未知数解释为面积比，最后做坐标代回。这会把几何、代数和条件检查放在同一条推理链上。
